% !PRINTED 482
\weberpage{482}
%<<CONT
then
\weq{21}{f_2(\omega) = \frac{e^{-\frac{r\pi i}{24}}\sqrt2}{x}, \quad f_2\left(\frac{\omega}{n}\right) = e^{-\frac{r\pi i}{24}}x,}
and these values are to be substituted for $u_1$, $v_1$ in the system II., § 73,
in order to obtain an equation for $x$.
By setting the various admissible values for $r$, one obtains
from (16):
\[
\begin{aligned}
n &= 3, & m &= 6,\ 5,\ 2,\\
n &= 5, & m &= 10,\ 9,\ 6,\ 1,\\
n &= 7, & m &= 14,\ 13,\ 10,\ 5,\\
n &= 11, & m &= 22,\ 21,\ 18,\ 13,\ 6,\\
n &= 13, & m &= 26,\ 25,\ 22,\ 17,\ 10,\ 1,\\
n &= 17, & m &= 34,\ 33,\ 30,\ 25,\ 18,\ 9,\\
n &= 19, & m &= 38,\ 37,\ 34,\ 29,\ 22,\ 13,\ 2.
\end{aligned}
\]
As a simple example we choose $n = 7$ and obtain from
§ 73, II.:
\[ \frac{x^8}{4} + \frac{4}{x^8} = 7 + 4\sqrt2 \cos\frac{r\pi}{4}. \]
From this there results for $r = 0$:
\weq{22}{\frac{x^2}{\sqrt2} + \frac{\sqrt2}{x^2} = 1+\sqrt2, \quad x = f_1(\sqrt{-14}),}
for $r=1$ the already known
\weq{23}{f(\sqrt{-13})^4 = 3+\sqrt{13},}
for $r=2$:
\weq{24}{f_1(\sqrt{-10})^2 = \frac{1+\sqrt5}{\sqrt2}.}
As a second example we take further $n = 13$ and obtain:
\weq{25}{A_1 = \frac{\sqrt2}{x^2} - \frac{x^2}{\sqrt2}, \quad B_1 = 16\cos\frac{r\pi}{2};}
for $r = 0, 4$, hence $m = 26, 10$, there results from this $B_1 = 16$, and
consequently, by § 73:
\[ A_1^7 - 6A_1^5 + A_1^3 + 20A_1 + 16 = 0, \]
while for $r=2$, hence $m=22$, in this equation $A_1$ is to be changed into
$-A_1$.
But one easily finds the decomposition:
\weq{26}{A_1^7 - 6A_1^5 + A_1^3 + 20A_1 + 16 = (A_1+1)^2(A_1-2)^2(A_1^3+2A_1^2+A_1+4).}
